% from ega0-15-fr.tex
\oldpage[0\textsubscript{IV-1}]{12}
Soient $A$ un anneau, $N$ un $A$-module~; dans ce qui suit, nous noterons
$N[T_1,\ldots,T_n]$, pour abréger, le $A$-module gradué
$N\otimes_A A[T_1,\ldots,T_n]$ ($T_i$ indéterminées), où $N$ est gradué par la
graduation triviale et $A[T_1,\ldots,T_n]$ par le degré total.

Soient $M$ un $A$-module, $f_i$ ($1\leq i\leq n$) des éléments de $A$,
$\mathfrak{J}=\sum_{i=1}^n f_iA$ l'idéal de $A$ engendré par les $f_i$, et
munissons $M$ de la filtration $\mathfrak{J}$-préadique~; on définit alors un
\emph{homomorphisme gradué surjectif de $A$-modules gradués}, de degré $0$,
\[
  \varphi : (\gr_0(M))[T_1,\ldots,T_n]\longrightarrow\gr_\bullet(M)
  \tag{15.1.1.1}\label{0.15.1.1.1-fr}
\]
de la façon suivante~: si $\xi\in\gr_0(M)$ est la classe mod.
$\mathfrak{J}M$ d'un élément $x\in M$, au couple formé de $\xi$ et d'un
polynôme $P[T_1,\ldots,T_n]$ homogène de degré $k$, on fait correspondre la
classe mod. $\mathfrak{J}^{k+1}M$ de $P(f_1,\ldots,f_n)x$~; il est immédiat
que cette classe
